\documentclass[12pt,a4paper]{ltjsarticle}
\usepackage{amsmath,amssymb}
\usepackage{bm}
\usepackage[margin=25mm]{geometry}
\usepackage{hyperref}
\usepackage{braket}
\usepackage{booktabs}
\usepackage{array}
\usepackage{tcolorbox}
\tcbuselibrary{breakable,skins}

% 初学者向けポイントボックス
\newtcolorbox{point}[1][]{
  colback=blue!5, colframe=blue!40!black,
  fonttitle=\bfseries, title=ポイント #1,
  breakable
}
\newtcolorbox{formula}[1][]{
  colback=green!5, colframe=green!40!black,
  fonttitle=\bfseries, title=公式 #1,
  breakable
}

\title{\textbf{$\beta^-$ 崩壊の基礎と論文まとめ}\\[6pt]
\large Large-scale shell model study of $\beta^-$-decay properties\\
of $N=126$, 125 nuclei\\[4pt]
\normalsize Phys.\ Rev.\ C \textbf{109}, 064319 (2024) \quad [arXiv:2309.07903]\\
\small S.~Sharma, P.~C.~Srivastava, A.~Kumar, T.~Suzuki, C.~Yuan, N.~Shimizu}
\date{2024年}

\begin{document}
\maketitle
\tableofcontents
\newpage

%==========================================================
\section{原子核の基礎知識}
%==========================================================

\subsection{原子核の構成}

原子核は \textbf{陽子（proton）} と \textbf{中性子（neutron）} から構成される．
これらをまとめて \textbf{核子（nucleon）} と呼ぶ．

\begin{center}
\begin{tabular}{lccc}
\toprule
粒子 & 記号 & 電荷 & 質量 \\
\midrule
陽子 & $p$ & $+e$ & $938.3$ MeV/$c^2$ \\
中性子 & $n$ & $0$ & $939.6$ MeV/$c^2$ \\
電子 & $e^-$ & $-e$ & $0.511$ MeV/$c^2$ \\
\bottomrule
\end{tabular}
\end{center}

原子核 $^A_Z X$（元素記号 $X$，質量数 $A$，陽子数 $Z$，中性子数 $N = A-Z$）において，
\textbf{陽子数}が元素の種類を決め，\textbf{中性子数}が変わると \textbf{同位体（isotope）} になる．

\begin{point}
\textbf{魔法数（magic number）}：
陽子数または中性子数が $2, 8, 20, 28, 50, 82, 126$ のとき，
原子核は特に安定になる（閉殻構造）．
本論文の対象核 $N=126$ はちょうど中性子の魔法数であり，構造が特殊である．
\end{point}

\subsection{核力と核構造}

核子同士は \textbf{核力（short-range）} で結びついている．
その基本的な特徴は：
\begin{itemize}
  \item 到達距離が短い（$\sim 1\text{--}2$ fm，1 fm $= 10^{-15}$ m）
  \item 陽子間・中性子間・陽子-中性子間でほぼ等しい（電荷独立性）
  \item スピンや軌道角運動量に依存する（スピン軌道力）
\end{itemize}

\subsection{スピンと角運動量}

核子はスピン $s = 1/2$ を持つフェルミ粒子である．
核の状態は角運動量量子数 $J$（全角運動量），パリティ $\pi = \pm 1$ で指定される（$J^\pi$ 表記）．

\begin{point}
\textbf{パウリの排他原理}：
同じ量子状態（量子数の組）に 2 つ以上の同種フェルミ粒子は入れない．
ただし \textbf{スピン上向き} と \textbf{スピン下向き} は別の状態なので，
各軌道に最大 2 個の核子が入れる（スピン自由度を含む）．
\end{point}

\subsection{同位旋（アイソスピン）}

陽子と中性子を統一的に扱うため \textbf{同位旋（isospin）} $t$ を導入する：
\begin{itemize}
  \item 陽子：$t_z = +1/2$
  \item 中性子：$t_z = -1/2$
  \item 降下演算子 $t_-$：$n \to p$（中性子を陽子に変換）$\Rightarrow$ $\beta^-$ 崩壊
  \item 上昇演算子 $t_+$：$p \to n$（陽子を中性子に変換）$\Rightarrow$ $\beta^+$ 崩壊
\end{itemize}

%==========================================================
\section{$\beta$ 崩壊の基礎}
%==========================================================

\subsection{$\beta$ 崩壊とは}

\textbf{$\beta$ 崩壊}は弱い相互作用（弱力）によって核子の種類が変わる崩壊過程である．

\begin{center}
\begin{tabular}{lll}
\toprule
種類 & 反応式 & 核種の変化 \\
\midrule
$\beta^-$ 崩壊 & $n \to p + e^- + \bar{\nu}_e$ & 中性子が陽子に変わる（$Z \to Z+1$） \\
$\beta^+$ 崩壊 & $p \to n + e^+ + \nu_e$ & 陽子が中性子に変わる（$Z \to Z-1$） \\
電子捕獲 (EC) & $p + e^- \to n + \nu_e$ & 陽子が中性子に変わる（$Z \to Z-1$） \\
\bottomrule
\end{tabular}
\end{center}

\begin{point}
本論文が扱う \textbf{$\beta^-$ 崩壊}は中性子過剰核（$N$ が多すぎる核）で起こる．
$N=126$ 付近の核は中性子が多く，$\beta^-$ 崩壊で安定に近づく方向に崩壊する．
\end{point}

\subsection{$Q$ 値（崩壊エネルギー）}

$\beta^-$ 崩壊が起こるためには，反応のエネルギー収支（$Q$ 値）が正である必要がある：
\begin{equation}
  Q_{\beta^-} = \bigl[M(Z, A) - M(Z+1, A)\bigr] c^2 > 0
\end{equation}
$M(Z, A)$：原子核の質量．$Q$ 値が大きいほど崩壊が速い（半減期が短い）傾向がある．

\subsection{フェルミ関数 $F(Z, W)$}

$\beta$ 崩壊で放出される電子は娘核のクーロン場の影響を受ける．
これを補正する係数が \textbf{フェルミ関数} $F(Z, W)$ である：
\begin{itemize}
  \item 重い核（$Z$ 大）ほどクーロン引力が強く，$F$ は大きくなる
  \item $W$：電子の全エネルギー（静止エネルギー $m_e c^2$ を含む）
\end{itemize}

\subsection{半減期・$\log ft$ 値}

\textbf{半減期} $t_{1/2}$ は原子核の個数が半分になるまでの時間：
\begin{equation}
  N(t) = N_0 \exp\!\left(-\frac{\ln 2}{t_{1/2}}\,t\right)
\end{equation}

$\beta$ 崩壊の強さを表す指標が \textbf{$\log ft$ 値}（log \textit{f} times \textit{t}）：
\begin{equation}
  \log ft = \log_{10}(f \cdot t_{1/2})
\end{equation}
\begin{itemize}
  \item $f$：位相空間因子（後述）
  \item $\log ft$ が小さいほど遷移が速い（許容遷移で $\log ft \sim 4$--$6$）
  \item $\log ft$ が大きいほど遷移が遅い（禁止遷移で $\log ft > 6$）
\end{itemize}

%==========================================================
\section{$\beta^-$ 崩壊の遷移分類}
%==========================================================

\subsection{選択則と遷移の種類}

$\beta$ 崩壊演算子を多極展開すると，角運動量変化 $\Delta J$ やパリティ変化 $\Delta\pi$ によって遷移が分類される：

\begin{center}
\begin{tabular}{lcccl}
\toprule
名称 & $\Delta J$ & $\Delta\pi$ & 典型 $\log ft$ & 備考 \\
\midrule
フェルミ（Fermi） & $0$ & $\pi\pi'=+1$ & $\sim 3.5$ & 同重体間の遷移 \\
ガモフ-テラー（GT） & $0, \pm 1$ & $\pi\pi'=+1$ & $4$--$6$ & スピンフリップ \\
1次禁止（FF） & $0, \pm 1, \pm 2$ & $\pi\pi'=-1$ & $6$--$9$ & 軌道角運動量変化 \\
2次禁止 & $\pm 2, \pm 3$ & $\pi\pi'=+1$ & $>10$ & さらに抑制 \\
\bottomrule
\end{tabular}
\end{center}

\begin{point}
本論文が扱う $N=126$ 核では \textbf{GT 遷移（許容）} と \textbf{1次禁止遷移（FF）} が
半減期を決める主要な成分である．フェルミ遷移はアイソスピン対称性で抑制される．
\end{point}

\subsection{ガモフ-テラー（GT）遷移の物理的意味}

GT 遷移では核子の \textbf{スピンが反転}しながら中性子が陽子に変わる：
\begin{equation}
  \hat{O}_{\rm GT} = g_A \sum_{k=1}^{A} \bm{\sigma}^{(k)} t_-^{(k)}
\end{equation}
\begin{itemize}
  \item $\bm{\sigma}$：スピン演算子（パウリ行列）
  \item $g_A = -1.2701(25)$：軸性ベクトル結合定数（実験値）
  \item $\sum_k$：全核子の和
\end{itemize}

\begin{point}
\textbf{パウリ・ブロッキング}：
ある軌道がすでに核子で占有されていると，そこへ遷移できない（パウリの排他原理）．
本論文では $\pi h_{11/2}$ 軌道の占有が $Z$ によって変わり，
GT 強度に直接影響することが示された．
\end{point}

\subsection{1次禁止（First-Forbidden; FF）遷移の物理的意味}

FF 遷移は，\textbf{軌道角運動量が 1 変化}する遷移で，
選択則のパリティが反転する（$\Delta\pi = -1$）：
\begin{equation}
  \hat{O}_{\rm FF} \sim i r \bigl[C_1 \times \bm{\sigma}\bigr]^J t_-, \quad J = 0, 1, 2
\end{equation}
$r$：核子の座標，$C_1$：1次球面調和テンソル．
演算子がより複雑で，GT より遷移確率が小さい（$\log ft$ が大きい）．

%==========================================================
\section{殻模型による $\beta$ 崩壊計算}
%==========================================================

\subsection{殻模型とは}

\textbf{殻模型（Shell Model）} は核子が平均場（ポテンシャル）の中で独立に運動するという
近似から出発し，残余相互作用を行列対角化で厳密に扱う手法：

\begin{center}
\fbox{
\begin{minipage}{0.85\textwidth}
\textbf{殻模型の手順}
\begin{enumerate}
  \item 平均場（ハーモニックオシレータ + スピン軌道力）で単粒子軌道を決める
  \item 閉殻（コア）外の核子に有効相互作用を与え，ハミルトニアン行列を構成する
  \item ハミルトニアン行列を対角化し，エネルギー固有状態（波動関数）を得る
  \item 得られた波動関数で物理量（遷移確率，電磁モーメント等）を計算する
\end{enumerate}
\end{minipage}
}
\end{center}

\subsection{研究の設定（本論文）}

\subsubsection{対象核種}
\begin{itemize}
  \item $N = 125, 126$（中性子魔法数 $N=126$ とその近傍），$Z = 52$--$79$
  \item $^{208}$Pb の「南」領域：$^{208}$Pb からみて陽子・中性子ともに少ない核種群
  \item r 過程元素合成（宇宙における重元素生成）に関わる核種
\end{itemize}

\subsubsection{使用コード：KSHELL}
核波動関数の生成には \textbf{KSHELL コード}（Shimizu ら開発）を使用する：
\begin{itemize}
  \item 大規模行列の対角化を並列計算で実行
  \item $N \leq 126$：模型空間の切り捨てなし（exact diagonalization）
  \item $N > 126$（コア破り核）：2粒子-2孔（2p-2h）励起を許容
\end{itemize}

\subsubsection{模型空間と有効相互作用}
\begin{center}
\begin{tabular}{ll}
\toprule
粒子種 & 軌道 \\
\midrule
陽子（$Z=50$--$82$） & $0g_{7/2}$, $1d_{5/2}$, $1d_{3/2}$, $2s_{1/2}$, $0h_{11/2}$ \\
中性子（$N=82$--$126$） & $0h_{9/2}$, $1f_{7/2}$, $1f_{5/2}$, $2p_{3/2}$, $2p_{1/2}$, $0i_{13/2}$ \\
\bottomrule
\end{tabular}
\end{center}
有効相互作用：KHHE（陽子間・中性子間）+ VM\_U+LS（陽子--中性子間）

%==========================================================
\section{$\beta^-$ 崩壊の計算公式}
%==========================================================

\subsection{GT 遷移強度 $B$(GT)}

\begin{formula}
\begin{equation}
  B({\rm GT}) = \frac{g_A^2}{2J_i + 1}
  \left| \Bra{f} \left\| \sum_k \bm{\sigma}^{(k)} t_-^{(k)} \right\| \Ket{i} \right|^2
\end{equation}
$J_i$：初期状態のスピン，$\Bra{f}\|\cdots\|\Ket{i}$：縮約行列要素（Wigner-Eckart の定理で計算）
\end{formula}

\subsubsection{クエンチング因子}
殻模型計算は GT 強度を過大評価する（短距離相関，$\Delta$ 粒子効果等が原因）．
実験値に合わせるため \textbf{クエンチング（quenching）因子} $q$ を導入：
\begin{equation}
  g_A^{\rm eff} = q \cdot g_A, \qquad q \simeq 0.74 \text{（本論文での決定値）}
\end{equation}
$q$ は実験の $\log ft$ との $\chi^2$ 最小化によって決定する．

\subsection{FF 遷移のシェイプファクター}

FF 遷移では，放出電子のエネルギー $W$ に依存するシェイプファクター $C(W)$ が現れる：
\begin{formula}
\begin{equation}
  C(W) = k_0 + k_1 W + k_{-1} W^{-1} + k_2 W^2
\end{equation}
\begin{itemize}
  \item 係数 $k_i$：ランク 0, 1, 2 の FF 行列要素と各クエンチング因子の組合せ
  \item クエンチング：$q_w = 0.41$，$q_x = 0.68$，$q_z = 0.76$（ランクごとに異なる）
\end{itemize}
\end{formula}

実験との比較には \textbf{平均シェイプファクター} を使用する：
\begin{equation}
  \langle C(W) \rangle = \frac{f}{f_0}
  \qquad
  \left(f_0 = \int_1^{W_0} F(Z,W)\, p\, W\, (W_0-W)^2\, dW\right)
\end{equation}

\subsection{位相空間因子 $f$}

\begin{formula}
\begin{equation}
  f = \int_1^{W_0} C(W)\, F(Z, W)\, p\, W\, (W_0 - W)^2\, dW
\end{equation}
\begin{itemize}
  \item $W_0 = Q_{\beta^-}/(m_e c^2) + 1$：最大電子エネルギー
  \item $p = \sqrt{W^2 - 1}$：電子の運動量（$m_e c$ 単位）
  \item $F(Z, W)$：フェルミ関数（娘核によるクーロン補正）
\end{itemize}
\end{formula}

GT 遷移の場合，$C(W) = B(\rm GT)$（エネルギー非依存）なので：
\begin{equation}
  f_{\rm GT} = B({\rm GT}) \cdot f_0
\end{equation}

\subsection{半減期と $\log ft$}

\begin{formula}
\begin{equation}
  t_{1/2} = \frac{\kappa}{f}, \qquad \kappa = 6289 \;\text{s}
\end{equation}
\begin{equation}
  \log ft = \log_{10}(f \cdot t_{1/2}^{\rm exp})
\end{equation}
GT 遷移では：
$\log ft_{\rm GT} = \log_{10}(\kappa / B({\rm GT}))$
\end{formula}

%==========================================================
\section{計算の全体手順}
%==========================================================

\begin{center}
\fbox{
\begin{minipage}{0.90\textwidth}
\textbf{$\beta^-$ 崩壊計算のフロー}
\begin{enumerate}
  \item \textbf{殻模型対角化}（KSHELLコード）\\
        $\Rightarrow$ 初期状態 $|i, J_i^\pi\rangle$，終状態 $|f, J_f^\pi\rangle$ の波動関数を取得
  \item \textbf{遷移行列要素の計算}\\
        GT：$\Bra{f}\|\sigma t_-\|\Ket{i}$\\
        FF：$\Bra{f}\|irC_1\!\times\!\sigma\, t_-\|\Ket{i}$（ランク 0, 1, 2）
  \item \textbf{クエンチング因子の適用}\\
        $g_A \to q \cdot g_A$（GT），$q_w, q_x, q_z$（FF）を各行列要素に乗じる
  \item \textbf{シェイプファクターの構成}\\
        GT：$C(W) = B(\rm GT)$（定数）\\
        FF：$C(W) = k_0 + k_1 W + k_{-1}/W + k_2 W^2$（エネルギー依存）
  \item \textbf{位相空間積分の数値計算}\\
        $Q_\beta$ 値を質量表から取得し，$f = \int_1^{W_0} C(W) F p W (W_0-W)^2 dW$ を計算
  \item \textbf{半減期・$\log ft$ 値の算出と実験値との比較}\\
        $t_{1/2} = 6289\,\text{s}/f$，$\log ft = \log_{10}(f \cdot t_{1/2}^{\rm exp})$
\end{enumerate}
\end{minipage}
}
\end{center}

%==========================================================
\section{論文の主な結果}
%==========================================================

\subsection{GT 強度分布の特徴}

\begin{itemize}
  \item GT 強度は \textbf{低励起エネルギー側（$E_x \lesssim 3$ MeV）に顕著なピーク}を持つ
  \item 支配的な遷移：$\nu\,0h_{9/2} \to \pi\,0h_{11/2}$（中性子が $h_{9/2}$ から陽子の $h_{11/2}$ へ）
  \item 陽子数が $Z < 82$ になると $\pi h_{11/2}$ の占有が減少
        $\Rightarrow$ \textbf{パウリ・ブロッキングが弱まり} GT ピークが強化される
\end{itemize}

\subsection{FF 遷移の役割}

\begin{itemize}
  \item $Z \sim 82$ 付近では FF 遷移が半減期に無視できない寄与をする
  \item 陽子数が減る（$Z \ll 82$）につれ，FF の寄与は相対的に小さくなる
\end{itemize}

\subsection{実験値との比較}

実験の $\log ft$ 値・平均シェイプファクターと広範に比較し，良好な一致を確認：

\begin{center}
\begin{tabular}{lll}
\toprule
核種 & 遷移 & 結果 \\
\midrule
$^{199}$Pt $\to$ $^{199}$Au & GT + FF & 実験値と良好に一致 \\
$^{199}$Au $\to$ $^{199}$Hg & GT + FF & 実験値と良好に一致 \\
$^{207}$Hg $\to$ $^{207}$Tl & GT + FF & 実験値と良好に一致 \\
$^{208}$Tl $\to$ $^{208}$Pb & GT + FF & 実験値と良好に一致 \\
\bottomrule
\end{tabular}
\end{center}

複数核種のスピン・パリティ（$J^\pi$）を理論的に新規提案した．

\subsection{r 過程元素合成への意義}

\begin{point}
\textbf{r 過程（rapid neutron-capture process）}：
超新星爆発や中性子星合体で起こる急速な中性子捕獲過程．
金，白金，ウランなどの重元素がこの過程で生成される．\\[4pt]
$N=126$ 待機核の正確な半減期データが r 過程計算の精度を決める．
本論文の結果は \textbf{信頼性の高い理論的半減期} を提供し，
元素合成シミュレーションの重要な入力データとなる．
\end{point}

%==========================================================
\section{HF.f への $\beta$ 崩壊実装の指針}
%==========================================================

\texttt{HF.f}（Skyrme TDHF コード）で収束した HF 基底状態の波動関数を用いて
$\beta$ 崩壊特性を計算する場合の実装方針（\texttt{Beta-Hatree.md} 参照）：

\begin{enumerate}
  \item \textbf{HF 波動関数の利用}\\
        虚時間法 (\texttt{DTexp}) で収束した一粒子軌道
        $\psi_n^{(p)}(\bm{r})$（中性子第 $p$ 軌道），
        $\psi_p^{(q)}(\bm{r})$（陽子第 $q$ 軌道）を使用する．

  \item \textbf{GT 行列要素の計算（実空間格子上）}\\
        格子点 $i$（座標 \texttt{xLval(i), yLval(i), zLval(i)}，ヤコビアン \texttt{Jacob(i)}）で：
        \begin{equation}
          M_{\rm GT}^{(\mu)} = \sum_i
          \bigl[\psi_p^{(q)*}(\bm{r}_i)\, \sigma_\mu\, \psi_n^{(p)}(\bm{r}_i)\bigr]
          \cdot \texttt{Jacob(i)} \cdot H^3
        \end{equation}
        スピン演算子 $\sigma_\mu$ は 2 成分スピノル $(\psi_{\uparrow}, \psi_{\downarrow})$ に作用する：
        $\sigma_z$：$\psi_1 \to \psi_1$，$\psi_2 \to -\psi_2$；
        $\sigma_\pm$：$\psi_1 \leftrightarrow \psi_2$

  \item \textbf{FF 行列要素の計算}\\
        \begin{equation}
          M_{\rm FF}^{J} = \sum_i r_i\,
          \bigl[\psi_p^{(q)*}(\bm{r}_i)\,
          [C_1(\hat{r}_i) \times \sigma]^J\,
          \psi_n^{(p)}(\bm{r}_i)\bigr]
          \cdot \texttt{Jacob(i)} \cdot H^3
        \end{equation}
        $r_i = \sqrt{x_i^2 + y_i^2 + z_i^2}$，$\hat{r}_i = (x_i, y_i, z_i)/r_i$

  \item \textbf{$Q$ 値と位相空間積分}\\
        $Q_{\beta^-}$ 値を質量表（AME2020 等）から外部入力として与え，
        フェルミ関数 $F(Z, W)$ を含む位相空間積分を数値計算する．

  \item \textbf{半減期の出力}\\
        $t_{1/2} = 6289\,\text{s} / f$ を計算し，\texttt{fort.XX} 等に書き出す．
\end{enumerate}

%==========================================================
\section{まとめ}
%==========================================================

\begin{center}
\begin{tabular}{>{\bfseries}p{4cm}p{9.5cm}}
\toprule
項目 & 内容 \\
\midrule
対象核種 & $N=125,126$，$Z=52$--$79$ の中性子過剰核 \\
崩壊過程 & $\beta^-$ 崩壊：$n \to p + e^- + \bar\nu_e$ \\
計算手法 & 大規模殻模型（KSHELL，切り捨てなし・2p-2h） \\
GT 演算子 & $g_A \sum \sigma t_-$，クエンチング $q \approx 0.74$ \\
FF 演算子 & $irC_1 \times \sigma$（ランク 0, 1, 2），各クエンチング因子 \\
位相空間 & $f = \int C(W) F(Z,W)\, p\, W\, (W_0-W)^2\, dW$ \\
半減期 & $t_{1/2} = 6289\,\text{s} / f$ \\
主な結果 & GT ピーク強度のパウリブロッキング依存性，FF の役割，実験値との一致 \\
物理的意義 & r 過程元素合成シミュレーションへの信頼性ある半減期データの提供 \\
\bottomrule
\end{tabular}
\end{center}

\vspace{1.5em}
\noindent\textbf{参考文献}\\
S.~Sharma et al., Phys.\ Rev.\ C \textbf{109}, 064319 (2024).
preprint: arXiv:2309.07903 [nucl-th].

\end{document}
